\begin{enumerate}[label=\textbf{\ansnum.\arabic*}]
  \item $\sigma^2=0.25$: $e=0.1$: $0.961$; $0.25$: $0.779$; $0.5$: $0.368$; $1.0$: $0.018$. The reward is $0.5$ at $e=\sigma\sqrt{\ln2}=0.5\cdot0.8326=0.416$\,m/s.
  \item $\sigma=0.3/\sqrt{\ln2}=0.360$. At $0.1$\,m/s: $\exp(-0.01/0.1296)=\exp(-0.077)=0.926$.
  \item First case: error $0.01+0.01+0.0025=0.0225$, term $\exp(-0.0225/0.25)=\exp(-0.09)=0.914$, contribution $1.83$. Second case: $(0.2-0.25)^2+(0-0.1)^2=0.0025+0.01=0.0125$, plus $v_z^2=0.0004$ gives $0.0129$; term $\exp(-0.0516)=0.950$, contribution $1.90$.
  \item $0.5\cdot0.4=0.2$\,rad/s. Reward $\exp(-0.16/0.5)=\exp(-0.32)=0.726$. For $1.5$\,rad: $0.75$ is clipped to $0.5$\,rad/s, the maximum turning rate.
  \item The rate is clipped at $0.5$\,rad/s while the error is above $1$\,rad: from $1.2$ to $1.0$ takes $0.2/0.5=0.4$\,s. Then $e(t)=1.0\,e^{-0.5\,t'}$ and $e<0.05$ needs $t'=2\ln20=5.99$\,s. In total about $6.4$\,s; a numerical loop with $dt=0.02$ gave $6.4$\,s. The time constant of the heading control is $2$\,s.
  \item When standing, the legs should not move and a tight tolerance holds the default pose. When running, large leg movements are needed; a tight tolerance would punish exactly the movements that produce the speed, and the robot would learn not to run (or to run with stiff legs and fall).
  \item Ankle roll: $\pm0.45$\,rad ($\pm26^{\circ}$) around the default; knee: $\pm1.0$\,rad ($\pm57^{\circ}$). A small angle error at the ankle moves the whole foot and the support polygon (Chapter~\ref{ch:ctl2}), so a limited target range keeps the policy from violent ankle commands.
  \item The policy that PPO finds is the one that minimises the penalties: stand still, or fall as soon as possible if the reward is negative and the termination ends the episode (a negative-return episode is better short). In the logs: the episode length falls, the velocity reward stays near zero, and the penalty terms dominate. The cure: reduce the penalty weights, or add a positive term (alive bonus) \emph{after} checking that the tracking terms work.
  \item Sketch: for each stage $k=0..9$ set $v_{\max}=0.2+0.8k/9$, draw $n$ commands with $p=0.2$ standing and $v_x\sim U(-v_{\max},v_{\max})$ otherwise, and print the share of standing time and the maximum of $|v_x|$. The expected standing share is $0.2$, weighted by the duration of the commands (which are uniform $3$--$8$\,s and independent of the type); the maximum of $|v_x|$ is close to $v_{\max}$ for many samples. (A sketch; I did not run it.)
  \item An example answer: actor: angular velocity, projected gravity, joint positions and velocities, last action, command (no linear velocity if your estimator is poor); critic: the same without noise plus the foot contact forces and the true linear velocity. Action scales as in the table of the K1 chapter. Rewards: the table of this chapter with the heading term and the pose term, plus foot air time with a small weight. Randomisation: pushes, friction, encoder bias, centre of mass, and later link masses and the action delay. Terminations: time-out, the trunk too low or too tilted. What to measure first on the real robot: the delay between a command and the motion of the joint and the torque limit, because they decide the randomisation range.
\end{enumerate}
